%------------------------------------------------------------------------------%


% 01
%------------------------------------------------------------------------------%
\begin{question}
  Descreva e esboce as curvas de nível da função \(f(x, y) = 100 - x^2 - y^2\)
\end{question}

\begin{resolution}{Solução}
  \onslide<+->{Para cada valor $c\leq 100$}

  \begin{align*}
    \onslide<+->{f(x, y) & = c \\}
    \onslide<+->{100 - x^2 - y^2 = c \\}
    \onslide<+->{- x^2 - y^2 = c - 100 \\}
    \onslide<+->{x^2 + y^2 = 100 - c}
  \end{align*}

  \onslide<+->{As curvas de nível são circunferências centradas nas origem de
    raio $\sqrt{100-c}$}
\end{resolution}

% 02
%------------------------------------------------------------------------------%
\begin{question}
  Descreva as superfícies de nível da função \(f(x, y, z) = \sqrt{x^2 + y^2 + z^2}\)
\end{question}

\begin{resolution}{Solução}
  \onslide<+->{Para cada valor $c \geq 0$}

  \begin{align*}
    \onslide<+->{f(x, y, z) & = c \\}
    \onslide<+->{\sqrt{x^2 + y^2 + z^2} & = c \\}
    \onslide<+->{x^2 + y^2 + z^2 & = c^2}
  \end{align*}

  \onslide<+->{As superfícies de nível são esferas centradas na origem de raio $c$}
\end{resolution}

% 03
%------------------------------------------------------------------------------%
\begin{question}
  Descreva e esboce as curvas de nível da função \(f(x, y) = xy\)
\end{question}

\begin{resolution}{Solução}
  \onslide<+->{Para cada valor $c \neq 0$}

  \begin{align*}
    \onslide<+->{f(x, y) & = c \\}
    \onslide<+->{xy & = c \\}
    \onslide<+->{y & = \frac{c}{x}}
  \end{align*}

  \onslide<+->{Para $c=0$, ou $x=0$ ou $y=0$}
\end{resolution}

% 04
%------------------------------------------------------------------------------%
\begin{question}
  Descreva e esboce as curvas de nível da função
  \(f(x, y) = \ln\left(9-x^2-y^2\right)\)
\end{question}

\begin{resolution}{Solução}
  \onslide<+->{Domínio de $f$}
  \begin{align*}
    \onslide<+->{9-x^2-y^2 & > 0         \\}
    \onslide<+->{9         & > x^2 + y^2 \\}
    \onslide<+->{x^2 + y^2 & < 3^2}
  \end{align*}
  \onslide<+->{Disco aberto centrado na origem de raio 3}

  \onslide<+->{Imagem:}
  \onslide<+->{a função assume valores no intervalo \(\big(-\infty, \ln(9)\,\big]\)}
\end{resolution}

\begin{resolution}{Solução}
  \onslide<+->{Para cada valor $c \leq \ln(9)$}

  \begin{align*}
    \onslide<+->{f(x, y)                   & = c       \\}
    \onslide<+->{\ln\left(9-x^2-y^2\right) & = c       \\}
    \onslide<+->{9-x^2-y^2                 & = e^c     \\}
    \onslide<+->{-x^2-y^2                  & = e^c - 9 \\}
    \onslide<+->{x^2+y^2                   & = 9 - e^c}
  \end{align*}

  \onslide<+->{As curvas de nível são circunferências centrada nas origem de
    raio $\sqrt{9 - e^c}$}
\end{resolution}

% 05
%------------------------------------------------------------------------------%
\begin{resolution}{Solução}
  Encontre a equação e esboce a curva de nível da função
  \[
    f(x, y) = \frac{2y-x}{x+y+1}
  \]
  que passa pelo ponto \((-1, 1)\)
\end{resolution}

\begin{resolution}{Solução}
  \begin{align*}
    \onslide<+->{c & = f(-1, 1)                              \\}
    \onslide<+->{c & = \frac{2y-x}{x+y+1} \evalat{(-1, 1)}{} \\}
    \onslide<+->{c & = \frac{2\times 1-(-1)}{(-1)+1+1}       \\}
    \onslide<+->{c & = \frac{2+1}{1}                         \\}
    \onslide<+->{c & = 3}
  \end{align*}
\end{resolution}

\begin{resolution}{Solução}
  \begin{align*}
    \onslide<+->{f(x, y)            & = 3        \\}
    \onslide<+->{\frac{2y-x}{x+y+1} & = 3        \\}
    \onslide<+->{2y-x               & = 3(x+y+1) \\}
    \onslide<+->{2y-x               & = 3x+3y+3  \\}
    \onslide<+->{2y-3y              & = 3x+x+3   \\}
    \onslide<+->{-y                 & = 4x+3     \\}
    \onslide<+->{y                  & = -4x -3}
  \end{align*}
\end{resolution}

% 06
%------------------------------------------------------------------------------%
\begin{resolution}{Solução}
  Encontre a equação e esboce a curva de nível da função
  \[
    f(x, y) = \sqrt{x^2-1}
  \]
  que passa pelo ponto \((1, 0)\)
\end{resolution}

\begin{resolution}{Solução}
Domínio
\begin{align*}
  \onslide<+->{x^2 -1  & \geq 0  \\}
  \onslide<+->{x^2     & \geq 1  \\}
  \onslide<+->{\abs{x} & \geq 1  \\}
  \onslide<+->{x       & \leq -1 \quad\text{ou}\quad x \geq 1}
\end{align*}
\end{resolution}

\begin{resolution}{Solução}
  \begin{align*}
    \onslide<+->{c & = f(1, 0)                           \\}
    \onslide<+->{c & = \sqrt{x^2-1} \; \evalat{(1, 0)}{} \\}
    \onslide<+->{c & = \sqrt{1-1}                        \\}
    \onslide<+->{c & = 0}
  \end{align*}
\end{resolution}

\begin{resolution}{Solução}
  \begin{align*}
    \onslide<+->{f(x, y)      & = 0 \\}
    \onslide<+->{\sqrt{x^2-1} & = 0 \\}
    \onslide<+->{x^2-1        & = 0 \\}
    \onslide<+->{x^2          & = 1 \\}
    \onslide<+->{\abs{x}      & = 1 \\}
    \onslide<+->{x            & = \pm 1}
  \end{align*}
\end{resolution}

% 07
%------------------------------------------------------------------------------%
\begin{resolution}{Solução}
  Encontre uma equação para a superfície de nível da função
  \[
    f(x, y, z) = \frac{x - y + z}{2x + y -z}
  \]
  que passa pelo ponto \((1, 0, -2)\)
\end{resolution}

\begin{resolution}{Solução}
  Domínio
  \begin{align*}
    \onslide<+->{2x + y -z  & \neq 0  \\}
    \onslide<+->{2x + y     & \neq z  \\}
    \onslide<+->{z          & \neq 2x + y}
  \end{align*}
\end{resolution}

\begin{resolution}{Solução}
  \begin{align*}
    \onslide<+->{c & = f(1, 0, -2)                                          \\}
    \onslide<+->{c & = \frac{x - y + z}{2x + y -z} \; \evalat{(1, 0, -2)}{} \\}
    \onslide<+->{c & = \frac{1 - 0 -2}{2\times 1 + 0 -(-2)}                 \\}
    \onslide<+->{c & = \frac{-1}{4}}
  \end{align*}
\end{resolution}

\begin{resolution}{Solução}
  \begin{align*}
    \onslide<+->{f(x, y, z)      & = -\frac{1}{4} \\}
    \onslide<+->{\frac{x - y + z}{2x + y -z} & = -\frac{1}{4} \\}
    \onslide<+->{-4(x - y + z)        & = 2x + y -z \\}
    \onslide<+->{-4x +4y -4z        & = 2x + y -z \\}
    \onslide<+->{-4z + z        & = 2x + y +4x -4y\\}
    \onslide<+->{-3z        & = 6x -3y\\}
    \onslide<+->{z        & = -2x + y}
  \end{align*}
\end{resolution}

% 08
%------------------------------------------------------------------------------%
\begin{resolution}{Solução}
Encontre a equação e esboce para a superfície de nível da função
\[
f(x, y, z) = \frac{x - y + z}{2x + y -z}
\]
que passa pelo ponto \((1, 0, -2)\)
\end{resolution}

\begin{resolution}{Solução}
Domínio
\begin{align*}
  \onslide<+->{2x + y -z  & \neq 0  \\}
  \onslide<+->{2x + y     & \neq z  \\}
  \onslide<+->{z          & \neq 2x + y}
\end{align*}
\end{resolution}

\begin{resolution}{Solução}
\begin{align*}
  \onslide<+->{c & = f(1, 0, -2)                                          \\}
  \onslide<+->{c & = \frac{x - y + z}{2x + y -z} \; \evalat{(1, 0, -2)}{} \\}
  \onslide<+->{c & = \frac{1 - 0 -2}{2\times 1 + 0 -(-2)}                 \\}
  \onslide<+->{c & = \frac{-1}{4}}
\end{align*}
\end{resolution}

\begin{resolution}{Solução}
\begin{align*}
  \onslide<+->{f(x, y, z)                  & = -\frac{1}{4} \\}
  \onslide<+->{\frac{x - y + z}{2x + y -z} & = -\frac{1}{4} \\}
  \onslide<+->{-4(x - y + z)               & = 2x + y -z \\}
  \onslide<+->{-4x +4y -4z                 & = 2x + y -z \\}
  \onslide<+->{-4z + z                     & = 2x + y +4x -4y\\}
  \onslide<+->{-3z                         & = 6x -3y\\}
  \onslide<+->{z                           & = -2x + y}
\end{align*}
\end{resolution}

% 09
%------------------------------------------------------------------------------%
\begin{resolution}{Solução}
  Encontre a equação e esboce para a superfície de nível da função
  \[
    f(x, y, z) = \sqrt{x-y} - \ln(z)
  \]
  que passa pelo ponto \((3, -1, 1)\)
\end{resolution}

\begin{resolution}{Solução}
  \onslide<+->{Domínio}
  \begin{align*}
    \onslide<+->{x-y & \geq 0  \\}
    \onslide<+->{x   & \geq y}
  \end{align*}
  \onslide<+->{e
  \[z > 0\]}

\end{resolution}

\begin{resolution}{Solução}
  \begin{align*}
    \onslide<+->{c & = f(3, -1, 1)                                  \\}
    \onslide<+->{c & = \sqrt{x-y} - \ln(z) \; \evalat{(3, -1, 1)}{} \\}
    \onslide<+->{c & = \sqrt{3-(-1)} - \ln(1)                       \\}
    \onslide<+->{c & = \sqrt{4}                                     \\}
    \onslide<+->{c & = 2}
  \end{align*}
\end{resolution}

\begin{resolution}{Solução}
  \begin{align*}
    \onslide<+->{f(x, y, z)          & = 2 \\}
    \onslide<+->{\sqrt{x-y} - \ln(z) & = 2 \\}
    \onslide<+->{\sqrt{x-y}          & = 2 + \ln(z) \\}
    \onslide<+->{-4x +4y -4z         & = 2x + y -z \\}
    \onslide<+->{-4z + z        & = 2x + y +4x -4y\\}
    \onslide<+->{-3z        & = 6x -3y\\}
    \onslide<+->{z        & = -2x + y}
  \end{align*}
\end{resolution}

%------------------------------------------------------------------------------%
